\subsection{Quantum Circuit Representation}

In the previous sections, we have been using visual representations of quantum gates and their effects on qubits. However, a quantum circuit can be expressed in a more formal way using a \textbf{single unitary matrix} that represents the entire operation of the circuit. This matrix acts on the global state of the system and is obtained by composing the individual gates according to the structure of the circuit. This is known as the \textbf{quantum circuit representation}.

\begin{definitionbox}[: Quantum Circuit Representation]
    A \textbf{quantum circuit} can be \textbf{represented as a unitary matrix} $U$ that acts on the global state vector of the qubits. The overall operation of the circuit is obtained by composing the unitary matrices corresponding to each layer of gates in the circuit.

    \highspace
    Mathematically, if a circuit consists of layers $G_1, G_2, \ldots, G_n$ with corresponding unitary matrices $U_1, U_2, \ldots, U_n$, then the overall unitary matrix $U$ representing the circuit is given by:
    \begin{equation*}
        U = U_n U_{n-1} \cdots U_2 U_1
    \end{equation*}
    Where the rightmost operator is applied first.

    \highspace
    Each $U_i$ is constructed by combining gates acting in parallel using tensor products, and by extending gates with identity operators on qubits where no operation is applied.
\end{definitionbox}

\noindent
In simple terms, a \textbf{quantum circuit is just a structured way to build a single unitary matrix acting on the whole system}. So:
\begin{equation*}
    \ket{\psi_{\text{out}}} = U \ket{\psi_{\text{in}}}
\end{equation*}
Where $\ket{\psi_{\text{in}}}$ is the initial state of the qubits, $\ket{\psi_{\text{out}}}$ is the final state after applying the circuit, and $U$ is the \textbf{product of all the gates in the circuit}.

\highspace
\begin{flushleft}
    \textcolor{Red2}{\faIcon{exclamation-triangle} \textbf{The 3 Fundamental Rules of Quantum Circuit Representation}}
\end{flushleft}
\begin{enumerate}
    \item \important{Parallel composition}. Gates acting on different qubits are combined using the tensor product:
    \begin{equation*}
        U = A \otimes B
    \end{equation*}

    \item \important{Sequential composition}. Gates applied in sequence are combined using matrix multiplication, evaluated from right to left:
    \begin{equation*}
        U = U_n \cdots U_2 U_1
    \end{equation*}
    So the first gate applied is the rightmost one in the product.

    \item \important{Identity extension}. When a gate does not act on a qubit, the identity operator must be used to extend it:
    \begin{equation*}
        U = I \otimes A
    \end{equation*}
    This keeps dimensions consistent when combining gates that act on different subsets of qubits.
\end{enumerate}

\begin{examplebox}[: Quantum Circuit Representation]
    Let's consider the following quantum circuit representation:
    \begin{equation*}
        \underbrace{(H \otimes I \otimes I)}_{4} \cdot \underbrace{(I \otimes \text{CNOT})}_{3} \cdot \underbrace{(X \otimes I \otimes H)}_{2} \cdot \underbrace{(\text{CNOT} \otimes I)}_{1}
    \end{equation*}
    First of all, the order of the gates is from right to left, so we start with the rightmost gate.
    \begin{center}
        \begin{quantikz}[slice all,slice style=blue]
            & \ctrl{1}  & \gate{X}  &           & \gate{H} & \\
            & \targ{}   &           & \ctrl{1}  &          & \\
            &           & \gate{H}  & \targ{}   &          &
        \end{quantikz}
    \end{center}
\end{examplebox}